\documentclass[aps,prl,reprint]{revtex4-2}
\usepackage[T1]{fontenc}
\usepackage[utf8]{inputenc}
\usepackage{amsmath,amssymb,mathtools,amsthm}
\usepackage[shortlabels]{enumitem}
\usepackage{booktabs}
\usepackage{longtable}
\usepackage{tabularx}
\usepackage{textcomp}
\usepackage{newunicodechar}
\usepackage[strings]{underscore}

\newunicodechar{→}{\ensuremath{\to}}
\newunicodechar{⇒}{\ensuremath{\Rightarrow}}
\newunicodechar{¬}{\ensuremath{\neg}}
\newunicodechar{≈}{\ensuremath{\approx}}
\newunicodechar{∘}{\ensuremath{\circ}}
\newunicodechar{⨟}{\ensuremath{\mathbin{\triangleright}}}
\newunicodechar{₅}{\textsubscript{5}}
\newunicodechar{²}{\textsuperscript{2}}
\newunicodechar{π}{\ensuremath{\pi}}
\newunicodechar{Π}{\ensuremath{\Pi}}
\newunicodechar{ω}{\ensuremath{\omega}}
\newunicodechar{̂}{\textasciicircum}
\newunicodechar{β}{\ensuremath{\beta}}
% --- Notation macros (readable, non-Agda) ---
\newcommand{\Sys}{\mathsf{Sys}}
\newcommand{\SK}{\mathsf{SK}}
\newcommand{\Time}{\mathsf{T}}
\newcommand{\RG}{\mathbb{N}}
\newcommand{\Diff}{\mathsf{Diff}}
\newcommand{\CSK}{\mathsf{C}}

\newcommand{\diag}{\operatorname{diag}}
\newcommand{\swap}{\operatorname{swap}}
\newcommand{\flip}{\operatorname{flip}}
\newcommand{\obs}{\operatorname{obs}}
\newcommand{\actT}{\operatorname{act}_t}
\newcommand{\actU}{\operatorname{act}_u}
\newcommand{\actomega}{\operatorname{act}_\omega}
\newcommand{\castR}[2]{\operatorname{cast}_R^{#1}\!\left(#2\right)}

\newcommand{\Ztwo}{\mathbb{Z}_2}

% --- Theorem environments ---
\theoremstyle{definition}
\newtheorem{definition}{Definition}[section]
\newtheorem{remark}{Remark}[section]
\theoremstyle{plain}
\newtheorem{lemma}{Lemma}[section]
\newtheorem{theorem}{Theorem}[section]
\newtheorem{proposition}{Proposition}[section]
\newtheorem{corollary}{Corollary}[section]

\begin{document}
\title{Two-time descent: bicomplex identities and a $\Ztwo$ obstruction (rigorous translation)}
\author{Dmytro Surdu}
\date{\today}
\maketitle

% This file is a rigorous informal (human-readable) proof translation of the
% It uses only standard mathematical notation.

\section*{Conventions}\label{sec:limited-scope-theorem-translation-clean-001}
All sets and maps are in the ordinary category of sets.
All groupoids are small.
Composition of morphisms is written left-to-right: for morphisms
$r\xrightarrow{p}s\xrightarrow{q}t$ we write $p\triangleright q:r\to t$ to mean
``do $p$ then $q$'', i.e. $q\circ p$ in the usual right-to-left convention.

\section{Assumptions (structural data)}
\label{sec:assumptions}
We fix the following abstract data.

\begin{enumerate}[label=\textbf{(A\arabic*)},leftmargin=*,itemsep=0.6em]
\item \textbf{Core SK data.}
A set of SK states $\SK$ together with a set of record states $\Sys$ and a map
\[\diag:\Sys\to\SK.\]
There is an involution (the SK swap)
\[\swap:\SK\to\SK,\qquad \swap\circ\swap=\mathrm{id}_{\SK},\qquad \swap\circ\diag=\diag.\]

\item \textbf{Two arrows (time and RG).}
Let $\Time$ be a group with identity $e_T$, and set $\RG:=\mathbb N$ with identity
$0$ and composition $+$.
We assume actions
\[
\actT:\Time\times\SK\to\SK,\qquad \actU:\RG\times\SK\to\SK
\]
satisfying the action laws
\[
\actT(e_T,x)=x,\qquad \actT(t_2t_1,x)=\actT(t_2,\actT(t_1,x)),\qquad
\actT(t^{-1},\actT(t,x))=x,
\]
\[
\actU(0,x)=x,\qquad \actU(u_2+u_1,x)=\actU(u_2,\actU(u_1,x)).
\]
Finally, weak commutativity holds on $\SK$:
\[
\actU(u,\actT(t,x))=\actT(t,\actU(u,x)).
\]

\item \textbf{Strict diagonalization of RG (record projection).}
For each $u\in\mathbb N$ with $u\ge 1$ there exists a map
\[
\Pi_{u}:\SK\to\Sys
\]
such that for all $x\in\SK$,
\[
\actU(u,x)=\diag(\Pi_{u}(x)).
\]
We also assume a record-level RG action
\[
R_u:\Sys\to\Sys
\]
with monoid laws
\[
R_0(r)=r,\qquad R_{u_2+u_1}(r)=R_{u_2}(R_{u_1}(r)),
\]
and intertwining on diagonal states:
\[
\Pi_{u}(\diag r)=R_{u}(r)\quad (u\ge 1),
\]
equivalently,
\[
\actU(u,\diag r)=\diag(R_u(r))\quad (u\ge 1).
\]

\item \textbf{RG-sector reachability (chosen scale).}
Fix $u_\star=1$ and assume a section
\[
\mathrm{sec}_\Pi:\Sys\to\SK
\]
such that
\[
\Pi_{1}(\mathrm{sec}_\Pi(r))=r.
\]

\item \textbf{No mixed anomaly (bisimplicial descent laws).}
There is a bisimplicial set $\CSK_{m,n}$ (for $m,n\ge 0$) equipped with
horizontal face maps $d^H_i:\CSK_{m+1,n}\to\CSK_{m,n}$ ($0\le i\le m+1$)
and vertical face maps $d^V_j:\CSK_{m,n+1}\to\CSK_{m,n}$ ($0\le j\le n+1$)
satisfying:
\begin{enumerate}[label=\textbf{(A5.\alph*)},leftmargin=*,itemsep=0.4em]
\item $d^H_i d^H_j = d^H_{j-1} d^H_i$ for $i<j$;
\item $d^V_i d^V_j = d^V_{j-1} d^V_i$ for $i<j$;
\item $d^H_i d^V_j = d^V_j d^H_i$ for all $i,j$.
\end{enumerate}

\item \textbf{Record reflection (diagonal cancellation).}
If two diagonal SK states agree, then the underlying records agree:
\[
  \diag(r) = \diag(r') \Rightarrow r = r'.
\]
Equivalently, the diagonal encoding of records has no redundancy.

\begin{remark}
\label{rem:limited-scope-theorem-translation-clean-001}
Assumption \textbf{(A5)} constrains only \emph{boundaries} (face maps) and thus
ensures the two-time square has a well-defined boundary. Assumption \textbf{(A6)}
is only a cancellation principle; record-level interchange is derived later from
\textbf{(A2)}, \textbf{(A3)}, \textbf{(A4)}, and \textbf{(A6)} (Lemma~\ref{lem:derived-interchange}). The
``before/after'' distinction is not encoded as two different equalities in
$\Sys$ (UIP may collapse them); it is encoded in the process presentation by a
binary tag $f\in\{\mathrm{diag},\mathrm{bridge}\}$ on square generators. The
meaning of this tag is fixed canonically by the Diff involution: $\mathrm{diag}$
acts as the identity and $\mathrm{bridge}$ acts as $\flip$ (Lemma~\ref{lem:canonical-fillers}).
Nontriviality of this tag is established later once a non-fixed point of $\flip$
is obtained (Corollary~\ref{cor:canonical-fillers-different}), which is a
consequence of the detector.
\end{remark}

\item \textbf{Real structure, universal pointer basis, and two-time mixing (analytic, used to derive a detector).}
There exists a complex Hilbert space $\mathcal H$ with an orthonormal basis
$\{\lvert a\rangle\}_{a\in\mathcal I}$ (the \emph{pointer basis}), and we identify
$\SK$ with a specified subset of trace-class operators on $\mathcal H$.

\medskip
\noindent
\textbf{Universality of the pointer basis (record algebra is fixed).}
Let $\mathsf{Diag}$ denote the set of trace-class operators that are diagonal in the pointer basis.
We assume that the record embedding is diagonal in this basis:
\[
  \diag(\Sys)\subseteq \mathsf{Diag}.
\]
Equivalently, for each $r\in\Sys$ there exist coefficients $\{p_a(r)\}_{a\in\mathcal I}$
such that $\diag(r)=\sum_a p_a(r)\lvert a\rangle\langle a\rvert$.
Moreover, all RG-diagonalized outputs land in the same diagonal subspace:
\[
  \actU(u,x)\in \mathsf{Diag}\qquad(\forall x\in\SK,\ \forall u\ge 1).
\]
This packages the scale-stability of records: the ``record basis'' does not depend on
when (or at which RG depth) diagonalization is applied~\cite{Zurek1981PointerBasis,JoosZeh1985Emergence,Zurek2003Decoherence}.

\medskip
\noindent
Let $\mathcal K$ denote the antiunitary complex conjugation in this pointer basis.

In this concrete instantiation we take $\Time=\mathbb R$ (additive group) and
use microscopic unitary evolution only to derive the detector witness. The
abstract $\actT$ in the assumptions is the effective evolution on RG-classical
states and is not assumed to preserve diagonality microscopically.

We assume:
\begin{enumerate}[label=\textbf{(A7.\alph*)},leftmargin=*,itemsep=0.4em]
\item (Swap as transpose/conjugation) the swap map is implemented by
\[\swap(\rho)=\mathcal K\,\rho^{\ast}\,\mathcal K^{-1},\]
so that in the pointer basis $\swap(\rho)=\rho^{\mathsf T}$.

\item (Time evolution) there is a nonzero bounded self-adjoint operator $H$ such that
\[\rho(t)=e^{-itH}\,\rho\,e^{itH}.\]

\item (Scale/RG mixing generator; weakmini core) there exists a dense subspace
$\mathcal D\subseteq\mathcal H$ such that $\mathcal D\subseteq\operatorname{Dom}(D)$,
$H\mathcal D\subseteq\mathcal D$, $D\mathcal D\subseteq\mathcal D$, $K\mathcal D=\mathcal D$,
and $KDK^{-1}=D$ on $\mathcal D$. The commutator is defined on $\mathcal D$ and satisfies
\begin{equation}
  (DH-HD)\psi=iH\psi \qquad (\forall\psi\in\mathcal D).
  \label{eq:mixing-comm}
\end{equation}

\item (Nondegeneracy) there exist indices $a\ne b$ and a record state $r\in\Sys$
whose diagonal density matrix $\rho:=\diag(r)=\sum_c p_c\lvert c\rangle\langle c\rvert$
satisfies $p_a\ne p_b$, and such that
\begin{equation}
  \Re(H_{ab})\ne 0.
  \label{eq:nondeg}
\end{equation}
\end{enumerate}

\end{enumerate}

\begin{remark}
\label{rem:limited-scope-theorem-translation-clean-002}
Assumption \textbf{(A7)} is the physics-motivated bridge used to \emph{derive} the
swap-odd detector required by the obstruction argument, rather than postulating it.
The purely combinatorial parts (bicomplex, $\omega$-cover, splitting obstruction)
use only \textbf{(A1)}--\textbf{(A6)}.
\end{remark}

\section{No real subgroup lemma and derived detector}
\label{sec:mixing-detector}

\subsection{Abstract algebraic lemma (formal core)}\label{subsec:limited-scope-theorem-translation-clean-001}

\begin{definition}[Abstract mixing interface]
\label{def:abstract-mixing}
Let $(G,+,0)$ be an abelian group. Fix endomorphisms
\[\mathrm{i}:G\to G,\qquad K:G\to G,\qquad \mathrm{ad}_D:G\to G.\]
Assume the identities
\begin{align}
  K(\mathrm{i}(x)) &= \mathrm{i}(-K(x)), &\text{(anti-linearity)}\\
  K(\mathrm{ad}_D(x)) &= \mathrm{ad}_D(K(x)), &\text{($K$-equivariance)}\\
  \mathrm{ad}_D(-x) &= -\mathrm{ad}_D(x), &\text{(oddness)}
\end{align}
and the cancellation properties
\begin{align}
  x=-x &\Rightarrow x=0, &\text{(no $2$-torsion)}\\
  \mathrm{i}(x)=0 &\Rightarrow x=0. &\text{(trivial kernel of $\mathrm{i}$)}
\end{align}
An element $H\in G$ satisfies the \emph{mixing relation} if $\mathrm{ad}_D(H)=\mathrm{i}(H)$.
It is \emph{$K$-odd} if $K(H)=-H$.

\end{definition}

\begin{remark}[Domain hygiene for the analytic instantiation]
In the Hilbert-space setting, we apply Definition~\ref{def:abstract-mixing} to the
additive group of bounded operators that preserve $\mathcal D$ and for which the
commutator with $D$ is defined on $\mathcal D$. Under \textbf{(A7.c)}, $\mathrm{ad}_D(H)$
is defined on $\mathcal D$ and equals $iH$ there.
\end{remark}

\begin{lemma}[No real subgroup]
\label{lem:no-real-subgroup}
Under Definition~\ref{def:abstract-mixing}, if an element $H\in G$ satisfies both
\[\mathrm{ad}_D(H)=\mathrm{i}(H)\qquad\text{and}\qquad K(H)=-H,\]
then $H=0$.

\end{lemma}

\begin{proof}
Using $K$-equivariance and the mixing relation,
\[\mathrm{ad}_D(KH)=K(\mathrm{ad}_D H)=K(\mathrm{i}H)=\mathrm{i}(-KH).\]
Substitute $KH=-H$ to obtain $\mathrm{ad}_D(-H)=\mathrm{i}(-(-H))=\mathrm{i}(H)$.
By oddness, $\mathrm{ad}_D(-H)=-\mathrm{ad}_D(H)$, hence
$-\mathrm{ad}_D(H)=\mathrm{i}(H)$. Replace $\mathrm{ad}_D(H)$ by $\mathrm{i}(H)$ to get
$-\mathrm{i}(H)=\mathrm{i}(H)$, i.e. $\mathrm{i}(H)=-\mathrm{i}(H)$.
No $2$-torsion gives $\mathrm{i}(H)=0$, and the trivial kernel of $\mathrm{i}$ gives $H=0$.
\end{proof}

\begin{remark}
Lemma~\ref{lem:no-real-subgroup} is the abstract form of the statement that, under
the weakmini core relation \eqref{eq:mixing-comm} on a $K$-invariant dense core and
a $K$-even scale generator $D$ on that core, one cannot have nontrivial dynamics
entirely in the ``real'' subgroup characterized by $KHK^{-1}=-H$.
\end{remark}

\subsection{Derived detector in the pointer basis}

\begin{lemma}[Swap-odd observable constant on the diagonal]
\label{lem:swap-odd-observable}
Fix distinct indices $a\ne b$ in the pointer basis and define the bounded self-adjoint
operator
\[O_{ab}:= i\big(\lvert a\rangle\langle b\rvert-\lvert b\rangle\langle a\rvert\big).\]
Define
\[\obs_{ab}(\rho):=\mathrm{Tr}(\rho\,O_{ab}).\]
Then for every \emph{self-adjoint} trace-class $\rho$:
\begin{enumerate}[label=\textbf{(\roman*)},leftmargin=*,itemsep=0.4em]
\item For every diagonal density matrix $\rho$ in the pointer basis,
$\obs_{ab}(\rho)=0$.
\item $\obs_{ab}(\swap\rho)=-\obs_{ab}(\rho)$.
\item Moreover $\obs_{ab}(\rho)=2\,\Im(\rho_{ab})$ where $\rho_{ab}=\langle a\rvert\rho\lvert b\rangle$.
\end{enumerate}

\end{lemma}

\begin{proof}
(i) If $\rho$ is diagonal then $\rho_{ab}=0$ for $a\ne b$, so $\mathrm{Tr}(\rho O_{ab})=0$.

(ii) In the pointer basis, $\swap(\rho)=\rho^{\mathsf T}$, hence
$(\swap\rho)_{ab}=\rho_{ba}=\overline{\rho_{ab}}$ for self-adjoint $\rho$.
Therefore $\Im((\swap\rho)_{ab})=-\Im(\rho_{ab})$, and (iii) below implies (ii).

(iii) Directly compute
\[\mathrm{Tr}(\rho O_{ab})= i(\rho_{ba}-\rho_{ab})=2\,\Im(\rho_{ab}),\]
using $\rho_{ba}=\overline{\rho_{ab}}$ for trace-class $\rho$ when $\rho$ is self-adjoint.
\end{proof}

\begin{lemma}[Nontriviality of the derived detector]
\label{lem:derived-detector-nontrivial}
Under Assumption \textbf{(A7)}, there exist indices $a\ne b$ and a (time-evolved)
state $\rho(t)$ such that
\[\obs_{ab}(\rho(t))\ne 0.\]
Consequently, the map $\obs:=\obs_{ab}$ is constant on diagonal record states and is
swap-odd, while not being identically zero.

\end{lemma}

\begin{proof}
Let $r\in\Sys$ and $\rho:=\diag(r)$ be as in \textbf{(A7.d)} and consider its
unitary time evolution $\rho(t)=e^{-itH}\rho e^{itH}$.
Differentiate at $t=0$:
\[\dot\rho(0)=-i[H,\rho].\]
The $(a,b)$-matrix element satisfies
\[\dot\rho_{ab}(0)=-iH_{ab}(p_b-p_a).\]
Write $H_{ab}=\Re(H_{ab})+i\Im(H_{ab})$. Then
\[\dot\rho_{ab}(0)=(p_b-p_a)\Im(H_{ab}) - i(p_b-p_a)\Re(H_{ab}).\]
Taking imaginary parts gives
\[\frac{d}{dt}\Im(\rho_{ab}(t))\Big\rvert_{t=0}=-(p_b-p_a)\Re(H_{ab}).\]
By \eqref{eq:nondeg}, the right-hand side is nonzero, hence for sufficiently small
$t$ we have $\Im(\rho_{ab}(t))\ne 0$, and by Lemma~\ref{lem:swap-odd-observable}(iii)
this implies $\obs_{ab}(\rho(t))\ne 0$.

Finally, Lemma~\ref{lem:swap-odd-observable}(i)--(ii) shows that $\obs_{ab}$ is
constant on diagonal record states and is swap-odd.
\end{proof}

\begin{remark}[Where Lemma~\ref{lem:no-real-subgroup} enters]
The nondegeneracy condition \eqref{eq:nondeg} requires a matrix element with nonzero
real part. The ``real subgroup loophole'' would allow choosing a pointer basis in
which $H$ is $K$-odd (hence purely imaginary off-diagonal entries and no real part).
Lemma~\ref{lem:no-real-subgroup} excludes that loophole when the weakmini core
relation \eqref{eq:mixing-comm} holds on a $K$-invariant dense core and $H\ne 0$.
Thus, under mixing, some basis element must have $\Re(H_{ab})\ne 0$ for some pair.
\emph{By the universality clause in \textbf{(A7)}, the diagonal record algebra is fixed across RG depth, so the indices $a,b$ appearing here label the same record outcomes in any procedure we compare.}
\end{remark}

\section{Two-time chains and bicomplex}
\label{sec:bicomplex}

\subsection{Free abelian chains}
For each $(m,n)$ set
\[C_{m,n}:=\mathbb Z[\CSK_{m,n}],\]
the free abelian group on the set $\CSK_{m,n}$. We write $[x]\in C_{m,n}$ for the
generator corresponding to $x\in\CSK_{m,n}$.

Each face map induces a homomorphism on chains by linear extension, still denoted
$d^H_i,d^V_j$, so $d^H_i([x])=[d^H_i x]$ and similarly for $d^V_j$.

\subsection{Two boundaries}\label{subsec:limited-scope-theorem-translation-clean-004}
Define homomorphisms
\begin{align}
  Q_t &: C_{m+1,n}\to C_{m,n},
  &Q_t([x]) &:= \sum_{i=0}^{m+1} (-1)^i\,[d^H_i x],
  \\
  Q_u &: C_{m,n+1}\to C_{m,n},
  &Q_u([y]) &:= \sum_{j=0}^{n+1} (-1)^{m+j}\,[d^V_j y].
\end{align}
The factor $(-1)^m$ in $Q_u$ is the Koszul twist by horizontal degree.

\begin{lemma}[Boundary-of-boundary expansion]
\label{lem:boundary-boundary}
Let $d_i$ be face maps of a simplicial object and define
$\partial=\sum_i (-1)^i d_i$. Then on generators,
\[
\partial^2
=\sum_{0\le i<j\le n}(-1)^{i+j}\big(d_i d_j - d_{j-1} d_i\big).
\]
\end{lemma}

\begin{proof}
Write
\[
\partial^2=\sum_{i,j}(-1)^{i+j}d_i d_j
\]
and split the sum into $i<j$ and $i\ge j$. For the $i\ge j$ part, set $i'=i+1$,
so $i'>j$ and $(-1)^{i+j}=-(-1)^{i'+j}$, giving
\[
\sum_{i\ge j}(-1)^{i+j}d_i d_j
=-\sum_{j<i'}(-1)^{i'+j}d_{i'-1}d_j.
\]
Rename $(i',j)$ as $(j,i)$ to obtain
\[
\sum_{i\ge j}(-1)^{i+j}d_i d_j
=-\sum_{i<j}(-1)^{i+j}d_{j-1}d_i,
\]
and combine with the $i<j$ part to get the claimed formula.
\end{proof}

\begin{theorem}[Bicomplex identities]
\label{thm:bicomplex}
For all $m,n\ge 0$ the maps satisfy
\[Q_t^2 = 0,\qquad Q_u^2=0,\qquad Q_tQ_u + Q_uQ_t = 0.\]

\end{theorem}

\begin{proof}
We prove each identity on generators.

\smallskip\noindent
\emph{(1) $Q_t^2=0$.}
Let $x\in\CSK_{m+2,n}$. Then
\[
Q_t^2([x])
=\sum_{0\le i<j\le m+2}(-1)^{i+j}\big[d^H_i d^H_j x - d^H_{j-1} d^H_i x\big]
\]
by Lemma~\ref{lem:boundary-boundary}. The simplicial identity
$d^H_i d^H_j = d^H_{j-1} d^H_i$ for $i<j$ makes every summand zero.

\smallskip\noindent
\emph{(2) $Q_u^2=0$.}
Let $y\in\CSK_{m,n+2}$. Then
\[
Q_u^2([y])
=\sum_{0\le j<k\le n+2}(-1)^{j+k}\big[d^V_j d^V_k y - d^V_{k-1} d^V_j y\big],
\]
again by Lemma~\ref{lem:boundary-boundary}; the vertical face law
$d^V_j d^V_k=d^V_{k-1}d^V_j$ kills each summand.

\smallskip\noindent
\emph{(3) Anticommutation.}
Let $z\in\CSK_{m+1,n+1}$. Expand
\[Q_tQ_u([z])=\sum_{j=0}^{n+1}\sum_{i=0}^{m+1}(-1)^{m+1+j+i}[d^H_i d^V_j z],\]
\[Q_uQ_t([z])=\sum_{i=0}^{m+1}\sum_{j=0}^{n+1}(-1)^{m+j+i}[d^V_j d^H_i z].\]
By mixed commutation $d^H_i d^V_j=d^V_j d^H_i$, the terms match, while the
coefficients differ by a global minus sign because $(-1)^{m+1+j+i}=-(-1)^{m+j+i}$.
Thus $Q_tQ_u([z])+Q_uQ_t([z])=0$.
\end{proof}

\section{Record process groupoid and the parity cocycle}
\label{sec:groupoid}
This section uses only Assumptions \textbf{(A1)}--\textbf{(A6)}.

\begin{definition}[Record-level maps]
\label{def:record-level-actions}
We write $R_u:\Sys\to\Sys$ for the record-level RG action from \textbf{(A3)}.
The record-level time action is derived from the RG-sector section \textbf{(A4)}:
\[
T_t(r):=\Pi_{1}(\actT(t,\mathrm{sec}_\Pi(r))).
\]
In particular, $R_0(r)=r$ and $R_{u_2+u_1}(r)=R_{u_2}(R_{u_1}(r))$, and for $u\ge 1$ we have
\[
\actU(u,\diag r)=\diag(R_u(r)).
\]

\end{definition}

\begin{lemma}[Derived diagonal equivariance on the RG sector]
\label{lem:derived-diag-equiv}
For all $t\in\Time$ and $r\in\Sys$,
\[
\actT(t,\diag r)=\diag(T_t(r)).
\]
\end{lemma}

\begin{proof}
By \textbf{(A4)} and diagonalization, $\diag r=\actU(1,\mathrm{sec}_\Pi(r))$.
Apply weak commutativity and the factorization of $\actU(1,\cdot)$ through $\diag$
to obtain the claim.
\end{proof}

\begin{lemma}[Derived record-time action laws]
\label{lem:derived-time-laws}
For all $r\in\Sys$ and $t,t_1,t_2\in\Time$,
\[
T_{e_T}(r)=r,\qquad
T_{t_2t_1}(r)=T_{t_2}(T_{t_1}(r)),\qquad
T_{t^{-1}}(T_t(r))=r.
\]
\end{lemma}

\begin{proof}
The identity law follows from $\actT(e_T,x)=x$ and the section property
$\Pi_{1}(\mathrm{sec}_\Pi(r))=r$.
For composition and inverse, observe that if $\Pi_{1}(x)=\Pi_{1}(y)$ then
$\Pi_{1}(\actT(t,x))=\Pi_{1}(\actT(t,y))$ by diagonalization, weak commutativity,
and record reflection. Apply this with
$x=\actT(t_1,\mathrm{sec}_\Pi(r))$ and $y=\mathrm{sec}_\Pi(T_{t_1}(r))$, and similarly
for inverses.
\end{proof}

\begin{lemma}[Derived record-level interchange]
\label{lem:derived-interchange}
Under Assumptions \textbf{(A2)}, \textbf{(A3)}, \textbf{(A4)}, and \textbf{(A6)}, for all
$t\in\Time$, $u\in\RG$, and $r\in\Sys$,
\[
  R_u(T_t(r)) = T_t(R_u(r)).
\]

\end{lemma}

\begin{proof}
If $u=0$, then $R_0=\mathrm{id}$ and the claim is immediate.
If $u\ge 1$, apply weak commutativity on $\diag(r)$, rewrite both sides using
Lemma~\ref{lem:derived-diag-equiv} and diagonalization, and cancel $\diag$ using
record reflection (A6).
\end{proof}

\begin{definition}[Typed syntax $\mathrm{Proc}_5$]
\label{def:Proc5}
For records $r,s\in\Sys$, define a set of morphisms $\mathrm{Proc}_5(r,s)$
inductively by the constructors:
\begin{enumerate}[label=\textbf{(G\arabic*)},leftmargin=*,itemsep=0.4em]
\item $\mathrm{id}_r : \mathrm{Proc}_5(r,r)$.
\item $\mathrm{time}(t,r): \mathrm{Proc}_5(r,T_t(r))$ for $t\in\Time$.
\item $\mathrm{rg}(u,r): \mathrm{Proc}_5(r,R_u(r))$ for $u\in\RG$.
\item $\mathrm{sq}(f;t,u,r): \mathrm{Proc}_5(R_u(T_t(r)),T_t(R_u(r)))$ for
$f\in\{\mathrm{diag},\mathrm{bridge}\}$.
\item Composition: if $p:\mathrm{Proc}_5(r,s)$ and $q:\mathrm{Proc}_5(s,t)$, then
$p\triangleright q:\mathrm{Proc}_5(r,t)$.
\item Inverse: if $p:\mathrm{Proc}_5(r,s)$ then $p^{-1}:\mathrm{Proc}_5(s,r)$.
\item Codomain cast: if $p:\mathrm{Proc}_5(r,s)$ and $e:s=s'$, then
$\castR{e}{p}:\mathrm{Proc}_5(r,s')$.
\end{enumerate}
\end{definition}

\begin{definition}[Presented equality $\approx_5$]
\label{def:limited-scope-theorem-translation-clean-001}
Let $\approx_5$ be the smallest relation on $\mathrm{Proc}_5(r,s)$ closed under:
\begin{enumerate}[label=\textbf{(R\arabic*)},leftmargin=*,itemsep=0.4em]
\item Equivalence rules (refl/sym/trans) and congruence under
composition, inverse, and cast.
\item Groupoid laws: associativity, left/right identity, inverse laws, and
$(p^{-1})^{-1}=p$, plus
\[
(p\triangleright q)^{-1}\approx_5 q^{-1}\triangleright p^{-1},\qquad
(\mathrm{id}_r)^{-1}\approx_5 \mathrm{id}_r.
\]
\item Cast law: $\castR{\mathrm{refl}}{p}\approx_5 p$.
\item Time laws (with explicit casts):
\[
\castR{T_{e_T}(r)=r}{\mathrm{time}(e_T,r)}\approx_5 \mathrm{id}_r,
\]
\[
\castR{T_{t_2t_1}(r)=T_{t_2}(T_{t_1}(r))}{\mathrm{time}(t_2t_1,r)}
\approx_5 \mathrm{time}(t_1,r)\triangleright \mathrm{time}(t_2,T_{t_1}(r)),
\]
\[
\castR{T_{t^{-1}}(T_t(r))=r}{
  \mathrm{time}(t,r)\triangleright \mathrm{time}(t^{-1},T_t(r))
}\approx_5 \mathrm{id}_r.
\]
\item RG laws (with explicit casts):
\[
\castR{R_0(r)=r}{\mathrm{rg}(0,r)}\approx_5 \mathrm{id}_r,
\]
\[
\castR{R_{u_2+u_1}(r)=R_{u_2}(R_{u_1}(r))}{\mathrm{rg}(u_2+u_1,r)}
\approx_5 \mathrm{rg}(u_1,r)\triangleright \mathrm{rg}(u_2,R_{u_1}(r)).
\]
\end{enumerate}
\emph{Crucially, there is \textbf{no} relation equating a composite containing
$\mathrm{sq}(f;t,u,r)$ with a composite that does not contain it.}
In particular, the two fillers $\mathrm{sq}(\mathrm{diag};t,u,r)$ and
$\mathrm{sq}(\mathrm{bridge};t,u,r)$ are not identified by any relation.
The operational meaning of the tag is fixed canonically by
Lemma~\ref{lem:canonical-fillers} (diag $\mapsto$ id, bridge $\mapsto$ flip on $\Diff$).
Nontriviality of the tag action is proved later from the detector
(Corollary~\ref{cor:canonical-fillers-different}).

\end{definition}

\begin{definition}[Presented groupoid $\mathcal G$]
\label{def:limited-scope-theorem-translation-clean-002}
Objects are records $r\in\Sys$, and morphisms are $\mathrm{Proc}_5(r,s)$ modulo
the relation $\approx_5$.
\end{definition}

\begin{definition}[Corner path and square-loop]
\label{def:sqloop}
For $f\in\{\mathrm{diag},\mathrm{bridge}\}$ define the canonical path to the
"upper-right corner"
\[
  \mathrm{toCorner}_f(t,u,r)
  :=\mathrm{time}(t,r)\triangleright\mathrm{rg}(u,T_t(r))\triangleright\mathrm{sq}(f;t,u,r),
\]
a morphism $r\to T_t(R_u(r))$.
The \emph{square-loop} is the loop at $r$ defined by
\[
  \mathrm{sqLoop}(t,u,r):=\mathrm{toCorner}_{\mathrm{diag}}(t,u,r)\triangleright
  \mathrm{toCorner}_{\mathrm{bridge}}(t,u,r)^{-1}.
\]
\end{definition}

\begin{definition}[Parity cocycle $\omega$]
\label{def:limited-scope-theorem-translation-clean-003}
Define $\omega$ on generators by
\[\omega(\mathrm{time})=0,\qquad \omega(\mathrm{rg})=0,\qquad
\omega(\mathrm{sq}(\mathrm{diag}))=0,\qquad \omega(\mathrm{sq}(\mathrm{bridge}))=1,\]
and extend by
\[
\omega(p\triangleright q)=\omega(p)+\omega(q)\pmod 2,\qquad
\omega(p^{-1})=\omega(p),\qquad
\omega(\castR{e}{p})=\omega(p).
\]
\end{definition}

\begin{lemma}[$\omega$ is well-defined and detects the square-loop]
\label{lem:omega}
The function $\omega$ respects all defining relations of $\mathcal G$ and thus
defines a functor $\omega:\mathcal G\to B\Ztwo$. Moreover,
\[\omega(\mathrm{sqLoop}(t,u,r))=1\qquad (t,u,r).\]

\end{lemma}

\begin{proof}
By definition, $\omega$ is additive under composition, invariant under inverses,
and ignores casts, so it respects the groupoid and cast laws. It also respects the
time and RG relations because those are generated entirely by the morphisms
$\mathrm{time}$ and $\mathrm{rg}$, on which $\omega$ is identically $0$.

Crucially, the presentation of $\mathcal G$ includes \emph{no} relation equating a
composite that contains a square filler $\mathrm{sq}(f;t,u,r)$ with one that does
not (Definition~\ref{def:Proc5}). Hence the assignments
$\omega(\mathrm{sq}(\mathrm{diag}))=0$ and $\omega(\mathrm{sq}(\mathrm{bridge}))=1$ are
consistent.

Finally, using Definition~\ref{def:sqloop},
\[
\omega(\mathrm{sqLoop}(t,u,r))
=\omega(\mathrm{toCorner}_{\mathrm{diag}}(t,u,r))
+\omega(\mathrm{toCorner}_{\mathrm{bridge}}(t,u,r))
=0+1=1,
\]
because $\omega$ vanishes on $\mathrm{time}$ and $\mathrm{rg}$ and differs only on the
chosen filler tag.
\end{proof}

\begin{corollary}[$\omega$ oddness on abelian coefficients]
\label{cor:omega-odd-abelian}
Let $G$ be an abelian group and define an action by
\[
\mathrm{act}^G_{\omega}(p)(x):=(-1)^{\omega(p)}x.
\]
Then for all $t,u,r$ and $x\in G$,
\[
\mathrm{act}^G_{\omega}(\mathrm{sqLoop}(t,u,r))(x)=-x.
\]

\end{corollary}

\begin{proof}
By Lemma~\ref{lem:omega}, $\omega(\mathrm{sqLoop})=1$, so $(-1)^{\omega(\mathrm{sqLoop})}=-1$.
\end{proof}

\section{The $\Ztwo$-cover and splitting obstruction}
\label{sec:cover}
This section uses only Assumptions \textbf{(A1)}--\textbf{(A6)}.

\begin{definition}[$\Ztwo$-cover of $\mathcal G$]
\label{def:limited-scope-theorem-translation-clean-004}
Objects are pairs $(r,b)$ with $r\in\Sys$ and $b\in\Ztwo$. A morphism
$(r,b)\to(s,c)$ is a morphism $p:r\to s$ in $\mathcal G$ such that
\[c=b+\omega(p)\pmod 2.\]
Composition and inverses are inherited from $\mathcal G$. The projection functor
$\pi:\widetilde{\mathcal G}\to\mathcal G$ forgets the parity.

\end{definition}

\begin{definition}[Splitting]
A \emph{splitting} is a functor $s:\mathcal G\to\widetilde{\mathcal G}$ with
$\pi\circ s=\mathrm{id}_{\mathcal G}$~\cite{Brown2006TopologyGroupoids}.

\end{definition}

\begin{lemma}[Uniqueness of lifts in the $\Ztwo$-cover]
\label{lem:unique-lift}
Fix objects $(r,b)$ and $(s,c)$ in $\widetilde{\mathcal G}$.
For a given base morphism $p:r\to s$ in $\mathcal G$, there exists a morphism
$(r,b)\to(s,c)$ lying over $p$ if and only if
\[
c=b+\omega(p)\pmod 2.
\]
When it exists, it is unique.

\end{lemma}

\begin{proof}
By definition of $\widetilde{\mathcal G}$, a morphism over $p$ is not extra
structure: it is exactly the datum of the base arrow $p$ together with the
parity constraint $c=b+\omega(p)$. Since $p$ and $(r,b)$ determine $c$ uniquely,
there can be at most one lift. Existence is exactly the satisfaction of the
constraint.
\end{proof}

\begin{lemma}[Splitting iff $\omega$ is a coboundary]
\label{lem:split-cob}
There exists a splitting if and only if there exists a function
$\sigma:\Sys\to\Ztwo$ such that for every morphism $p:r\to s$ in $\mathcal G$,
\[\omega(p)=\sigma(r)+\sigma(s)\pmod 2.\]

\end{lemma}

\begin{proof}
If $s$ is a splitting, define $\sigma(r)$ by $s(r)=(r,\sigma(r))$. The lifted arrow
$s(p)$ must satisfy $\sigma(s)=\sigma(r)+\omega(p)$ by the defining condition of
$\widetilde{\mathcal G}$.

Conversely, given $\sigma$ satisfying the displayed equation, lift each object to
$(r,\sigma(r))$. For each $p:r\to s$ in $\mathcal G$, define $s(p)$ to be the unique lift
$(r,\sigma(r))\to(s,\sigma(s))$ lying over $p$; this lift exists by the parity
equation $\sigma(s)=\sigma(r)+\omega(p)$ and is unique by Lemma~\ref{lem:unique-lift}.
Additivity of $\omega$ implies functoriality.
\end{proof}

\begin{corollary}[Odd loop obstructs splitting]
\label{cor:odd-loop}
If there exists a loop $p:r\to r$ with $\omega(p)=1$, then no splitting exists.

\end{corollary}

\begin{proof}
If a splitting existed, Lemma~\ref{lem:split-cob} would imply
$\omega(p)=\sigma(r)+\sigma(r)=0$, contradiction.
\end{proof}

\begin{theorem}[Non-splitting of the $\Ztwo$-cover]
\label{thm:non-splitting}
Under Assumptions \textbf{(A1)}--\textbf{(A6)}, if $\Sys$ is nonempty then the $\Ztwo$-cover
$\pi:\widetilde{\mathcal G}\to\mathcal G$ does not admit a splitting.
Equivalently, the cocycle $\omega$ is not a coboundary on $\mathcal G$.
(If $\Sys=\emptyset$, the claim is trivial.)

\end{theorem}

\begin{proof}
If $\Sys$ is empty there is nothing to prove. Otherwise fix $r\in\Sys$, take $t=e_T$
and $u=1$, and set $p=\mathrm{sqLoop}(t,u,r)$. If a splitting existed,
Lemma~\ref{lem:split-cob} would imply $\omega(p)=0$ for every loop.
But Lemma~\ref{lem:omega} gives $\omega(p)=1$, a contradiction.
\end{proof}

\begin{proposition}[Two-way corner-path witness of the splitting obstruction]
\label{prop:two-way-obstruction}
Fix any $t\in\Time$, $u\in\RG$, and $r\in\Sys$, and consider the two parallel morphisms
\[
p_{\mathrm{diag}}:=\mathrm{toCorner}_{\mathrm{diag}}(t,u,r),\qquad
p_{\mathrm{bridge}}:=\mathrm{toCorner}_{\mathrm{bridge}}(t,u,r)
\]
from $r$ to $T_t(R_u(r))$ in $\mathcal G$.
Then $\omega(p_{\mathrm{diag}})=0$ and $\omega(p_{\mathrm{bridge}})=1$.
Consequently, no splitting of $\widetilde{\mathcal G}\to\mathcal G$ exists (as in Theorem~\ref{thm:non-splitting});
equivalently, the obstruction is already witnessed at this two-way corner configuration.

\end{proposition}

\begin{proof}
By definition, $\omega$ vanishes on $\mathrm{time}$ and $\mathrm{rg}$ and equals $0$ on $\mathrm{sq}(\mathrm{diag})$
and $1$ on $\mathrm{sq}(\mathrm{bridge})$, hence the two values are $0$ and $1$.

If a splitting existed, Lemma~\ref{lem:split-cob} would give $\sigma:\Sys\to\Ztwo$ with
$\omega(p)=\sigma(\mathrm{src}(p))+\sigma(\mathrm{tgt}(p))$. Since $p_{\mathrm{diag}}$ and $p_{\mathrm{bridge}}$
have the same source and target, this would force $\omega(p_{\mathrm{diag}})=\omega(p_{\mathrm{bridge}})$, contradiction.
\end{proof}

\begin{corollary}[Three-checkpoint loop witness]
\label{cor:groundhog-loop}
Assume there exist a record state $r\in\Sys$, a time step $t\in\Time$, and an RG step
$u\in\RG$ such that both procedures close at the record level:
\[
R_u(T_t(r))=r
\qquad\text{and}\qquad
T_t(R_u(r))=r.
\]
(Equivalently, it suffices to assume one of these equalities once record-level
interchange $T_t\!\circ R_u = R_u\!\circ T_t$ has been derived.)

Define the intermediate record checkpoints
\[
A:=r,\qquad B:=T_t(r),\qquad C:=R_u(r).
\]
Then the two corner procedures of Proposition~\ref{prop:two-way-obstruction},
\[
p_{\mathrm{diag}},\ p_{\mathrm{bridge}}:\ A\longrightarrow T_t(R_u(A))=A,
\]
are based loops at $A$ whose $\Ztwo$-labels differ:
\[
\omega(p_{\mathrm{diag}})=0,\qquad \omega(p_{\mathrm{bridge}})=1.
\]
Equivalently, even when the start and end coincide (three checkpoints), the distinction
between ``forgetting after'' and ``forgetting before'' carries an intrinsic one-bit holonomy
that cannot be reduced to a state-label on records.
\end{corollary}

\begin{proof}
Immediate from Proposition~\ref{prop:two-way-obstruction} after substituting the
assumption $T_t(R_u(r))=r$, which turns both parallel morphisms into loops based at $r$.
\end{proof}

\begin{proposition}[Collapsed presentation splits]
\label{prop:collapsed-splits}
Let $\mathcal G^{\mathrm{coll}}$ be the quotient of $\mathcal G$ obtained by adding the relations
\[
\mathrm{sq}(\mathrm{diag};t,u,r)\approx_5 \mathrm{sq}(\mathrm{bridge};t,u,r)\qquad (t,u,r),
\]
equivalently, identifying the two corner paths for each square. Then any functor
$\mathcal G^{\mathrm{coll}}\to B\Ztwo$ is trivial, and the induced $\Ztwo$-cover splits.
\end{proposition}

\begin{proof}
The added relations force the images of the two square generators to agree, so the
assignment $\omega(\mathrm{sq}(\mathrm{diag}))=0$, $\omega(\mathrm{sq}(\mathrm{bridge}))=1$
cannot descend. The only parity functor on $\mathcal G^{\mathrm{coll}}$ is therefore the
trivial one, which is a coboundary (take $\sigma\equiv 0$), so the cover splits.
\end{proof}

\begin{remark}[Tag collapse vs obstruction]
\label{rem:limited-scope-theorem-translation-clean-006}
Theorem~\ref{thm:non-splitting} concerns the non-collapsed pair $(\mathcal G,\omega)$.
If one collapses the filler tags, one changes the input to $(\mathcal G^{\mathrm{coll}},
\omega_{\mathrm{triv}})$ as in Proposition~\ref{prop:collapsed-splits}. The semantic
package in Section~\ref{sec:nontrivial} shows that such a collapse is incompatible with
the intended interpretation on $\Diff$: the detector supplies a non-fixed point of $\flip$,
so the induced $\omega$-action is genuinely nontrivial.
\end{remark}

\section{Difference object, flip, and the $\omega$-action}
\label{sec:diff}
This section uses only Assumptions \textbf{(A1)}--\textbf{(A6)}.

\begin{definition}[Set-level cofiber]
\label{def:limited-scope-theorem-translation-clean-006}
For a map $f:A\to B$, define $\mathrm{Cofib}(f)$ to be the quotient of
$B\sqcup\{\ast\}$ by the relation that identifies all points $f(a)$ with the basepoint
$\ast$. Let $\iota:B\to\mathrm{Cofib}(f)$ denote the induced map.

Define
\[\Diff:=\mathrm{Cofib}(\diag),\qquad \iota:\SK\to\Diff.\]

\end{definition}

\begin{lemma}[Cofiber universal property]
\label{lem:cofiber-up}
Let $f:A\to B$ and let $X$ be a set. If $\phi:B\to X$ is constant on $\mathrm{im}(f)$,
i.e. $\phi(f(a))=x_0$ for some $x_0\in X$ and all $a\in A$, then there exists a unique
$\overline\phi:\mathrm{Cofib}(f)\to X$ with $\overline\phi(\iota(b))=\phi(b)$ and
$\overline\phi(\ast)=x_0$.

\end{lemma}

\begin{proof}
Define $\overline\phi$ on representatives by $\overline\phi(\iota(b)):=\phi(b)$ and
$\overline\phi(\ast):=x_0$. The hypothesis ensures this respects the identifying
relation, hence descends to the quotient; uniqueness is immediate.
\end{proof}

\begin{definition}[Flip involution]
The commuting square $\swap\circ\diag=\diag$ induces an endomorphism of the cofiber,
\[\flip:\Diff\to\Diff,\qquad \flip(\iota(x)):=\iota(\swap x),\qquad \flip(\ast):=\ast.\]
Then $\flip^2=\mathrm{id}_{\Diff}$.

\end{definition}

\begin{lemma}[Canonical two fillers]
\label{lem:canonical-fillers}
Define $\operatorname{toggle}(\mathrm{diag})=\mathrm{bridge}$ and $\operatorname{toggle}(\mathrm{bridge})=\mathrm{diag}$, and
define the filler action on $\Diff$ by
\[
\mathrm{actFill}(\mathrm{diag})=\mathrm{id}_{\Diff},\qquad
\mathrm{actFill}(\mathrm{bridge})=\flip.
\]
Formally, $\mathrm{actFill}$ is defined by composing the tag parity
$\mathrm{fillParity}:\{\mathrm{diag},\mathrm{bridge}\}\to\{\mathrm{false},\mathrm{true}\}$
with the Bool-action on $\Diff$ induced by $\flip$.
Then for all tags $f$ and all $x\in\Diff$,
\[
\mathrm{actFill}(\operatorname{toggle} f)(x)=\flip(\mathrm{actFill}(f)(x)),\qquad
\operatorname{toggle}(\operatorname{toggle} f)=f.
\]
In particular, the orbit generated by $\flip$ is contained in $\{\mathrm{id},\flip\}$
(so there are no further canonical choices; these two may coincide if $\flip=\mathrm{id}$).

\end{lemma}

\begin{definition}[$\omega$-action on $\Diff$]
\label{def:limited-scope-theorem-translation-clean-008}
The involution $\flip$ defines a $\Ztwo$-action by $0\cdot y:=y$ and $1\cdot y:=\flip(y)$.
For each morphism $p$ in $\mathcal G$, define
\[\actomega(p)(y):=\omega(p)\cdot y.\]

\end{definition}

\begin{lemma}[Square-loop acts by flip]
\label{lem:sqLoop-flip}
For all $t,u,r$ and $y\in\Diff$,
\[\actomega(\mathrm{sqLoop}(t,u,r))(y)=\flip(y).\]

\end{lemma}

\begin{proof}
By Lemma~\ref{lem:omega}, $\omega(\mathrm{sqLoop}(t,u,r))=1$, and by definition of the
$\Ztwo$-action, $1\cdot y=\flip(y)$.
\end{proof}

\section{Semantic nontriviality from the derived detector}
\label{sec:nontrivial}
From here we use Assumption \textbf{(A7)}.

\begin{lemma}[Derived detector gives a non-fixed point in $\Diff$]
\label{lem:nonfixed}
Let $\obs:=\obs_{ab}$ be the observable produced by Lemma~\ref{lem:derived-detector-nontrivial}.
Then there exists $y\in\Diff$ with $\flip(y)\ne y$.

\end{lemma}

\begin{proof}
By Lemma~\ref{lem:swap-odd-observable}(i), $\obs$ is constant on diagonal record
states, so Lemma~\ref{lem:cofiber-up} yields a unique factorization
$\overline\obs:\Diff\to\mathbb R$ with $\overline\obs(\iota(x))=\obs(x)$.
Choose $x\in\SK$ such that $\obs(\swap x)\ne \obs(x)$ (exists by
Lemma~\ref{lem:derived-detector-nontrivial} and swap-oddness). Set $y:=\iota(x)$.
If $\flip(y)=y$, then applying $\overline\obs$ gives
$\obs(\swap x)=\obs(x)$, contradiction. Hence $\flip(y)\ne y$.
\end{proof}

\begin{corollary}[Distinguishability of canonical fillers]
\label{cor:canonical-fillers-different}
Under Assumption \textbf{(A7)}, there exists $y\in\Diff$ such that
\[
\mathrm{actFill}(\mathrm{bridge})(y)\neq \mathrm{actFill}(\mathrm{diag})(y).
\]
\end{corollary}

\begin{proof}
By Lemma~\ref{lem:nonfixed} choose $y$ with $\flip(y)\ne y$. By
Lemma~\ref{lem:canonical-fillers}, $\mathrm{actFill}(\mathrm{diag})=\mathrm{id}$
and $\mathrm{actFill}(\mathrm{bridge})=\flip$, hence the two values differ at $y$.
\end{proof}

\begin{theorem}[Semantic nontriviality (detectable holonomy)]
\label{thm:semantic-nontriviality}
Under Assumption \textbf{(A7)}, the $\omega$-action on $\Diff$ is nontrivial:
there exists $y\in\Diff$ with $\flip(y)\neq y$, and the square-loop acts
nontrivially on $\Diff$. In particular, the canonical fillers are distinguishable.

\end{theorem}

\begin{proof}
Lemma~\ref{lem:nonfixed} gives a non-fixed point of $\flip$, and
Lemma~\ref{lem:sqLoop-flip} identifies the $\omega$-action of $\mathrm{sqLoop}$
with $\flip$. Distinguishability follows from Corollary~\ref{cor:canonical-fillers-different}.
\end{proof}

\begin{remark}[Why $\Diff$ is essential for semantics]
\label{rem:limited-scope-theorem-translation-clean-007}
The oddness $\omega(\mathrm{sqLoop})=1$ is already a cohomological obstruction to
splitting (Theorem~\ref{thm:non-splitting}); $\Diff$ is \emph{not} needed for that.
What $\Diff$ provides is a concrete $\Ztwo$-action in which $\omega$ is realized as
monodromy. The detector supplies a non-fixed point of $\flip$, showing that this
monodromy is genuinely nontrivial and the filler tags are operationally distinguishable.
\end{remark}

\begin{theorem}[Limited Scope Theorem]
\label{thm:limited-scope}
Under Assumptions \textbf{(A1)}--\textbf{(A7)}, the following hold:
\begin{enumerate}[label=\textbf{(\arabic*)},leftmargin=*,itemsep=0.4em]
\item \textbf{Cohomological package (uses \textbf{(A1)}--\textbf{(A6)}).}
The bicomplex identities $Q_t^2=0$, $Q_u^2=0$, and $Q_tQ_u+Q_uQ_t=0$
(Theorem~\ref{thm:bicomplex}) hold; the presented groupoid $\mathcal G$ and cocycle
$\omega$ are defined with $\omega(\mathrm{sqLoop})=1$ (Lemma~\ref{lem:omega});
the $\Ztwo$-cover splits iff $\omega$ is a coboundary (Lemma~\ref{lem:split-cob});
the cover does not split (Theorem~\ref{thm:non-splitting}); and for any abelian
group $G$ the induced $\omega$-action satisfies
$\mathrm{act}^G_{\omega}(\mathrm{sqLoop}(t,u,r))(x)=-x$
(Corollary~\ref{cor:omega-odd-abelian}).
\item \textbf{Semantic package (uses \textbf{(A7)}).}
The $\omega$-action on $\Diff$ is defined and $\mathrm{sqLoop}$ acts by $\flip$
(Lemma~\ref{lem:sqLoop-flip}); the detector gives a non-fixed point of $\flip$
and hence distinguishable canonical fillers (Theorem~\ref{thm:semantic-nontriviality}).
\end{enumerate}
\end{theorem}

\begin{proof}
Item (1) collects the cited cohomological results; item (2) is the semantic
nontriviality conclusion.
\end{proof}

\begin{thebibliography}{9}
\bibitem{Zurek1981PointerBasis}
W.~H.~Zurek, Phys.\ Rev.\ D \textbf{24}, 1516 (1981), doi:10.1103/PhysRevD.24.1516.

\bibitem{JoosZeh1985Emergence}
E.~Joos and H.~D.~Zeh, Z.\ Phys.\ B \textbf{59}, 223 (1985), doi:10.1007/BF01725541.

\bibitem{Zurek2003Decoherence}
W.~H.~Zurek, Rev.\ Mod.\ Phys.\ \textbf{75}, 715 (2003), doi:10.1103/RevModPhys.75.715.

\bibitem{Brown2006TopologyGroupoids}
R.~Brown, \emph{Topology and Groupoids} (www.groupoids.org, Deganwy, United Kingdom, 2006).
\end{thebibliography}

\end{document}
